\section{Introduction}
\label{sec:introduction}

Cloud-native architectures orchestrated by Kubernetes have fundamentally transformed enterprise software delivery by decomposing monolithic applications into polyglot, ephemerally scheduled microservices \cite{beyer_2016, majors_2022}. While this architectural shift maximizes deployment velocity and horizontal elasticity, it replaces deterministic in-memory function calls with stochastic network interactions across virtualized kernel namespaces (\texttt{cgroups}, \texttt{netns}, and Container Network Interfaces). Consequently, when performance degradations, resource contention, or partial network partitions occur, traditional perimeter monitoring and reactive host polling fail to reconstruct the internal execution state of the distributed system \cite{sridharan_2018, sigelman_2010}.

To guide organizations from opaque, reactive operations toward proactive fault localization, observability vendors, industry consortia, and recent academic studies have proposed various ``Observability Maturity Models'' \cite{grafana_omm_2024, dynatrace_omm_2024, wwt_omm_2025, patel_2025}. However, a systematic inspection grounded in a Rapid Review of the cloud-native literature \cite{cartaxo_2018} (screening 312 studies to synthesize the 10 primary Kubernetes reference works \cite{kosinska_2023, usman_2022, faseeha_2025, araujo_2024, hartikainen_2024, gomes_2025, han_2024, deng_2024, zanatta_23, wei_2025}) alongside Design Science Research (DSR) \cite{hevner_2004, becker_2009} and maturity model meta-theory \cite{debruin_2005, wendler_2012, lasrado_2015} exposes three fundamental scientific deficiencies across existing frameworks:
\begin{enumerate}[leftmargin=*, nosep, label=(\roman*)]
    \item \textbf{Vendor Lock-in and Telemetry Volume Bias:} Industry maturity models (e.g., Grafana Labs \cite{grafana_omm_2024}, Dynatrace \cite{dynatrace_omm_2024}, WWT \cite{wwt_omm_2025}) equate higher maturity levels with the ingestion of larger telemetry volumes or the acquisition of proprietary SaaS modules, lacking vendor-neutral architectural criteria (such as OpenTelemetry and W3C TraceContext compliance).
    \item \textbf{Probe Effect and FinOps Blindness:} Existing models assume a strictly linear benefit to deeper instrumentation, ignoring the \textit{Probe Effect} ($\Omega$)---the marginal CPU, memory, and network bandwidth overhead imposed by user-space SDKs, sidecars, and collectors on production workloads \cite{agarwal_2022, valbusa_2025}.
    \item \textbf{The Monomodal Additivity Fallacy:} Current frameworks implicitly assume that adopting individual telemetry pillars (e.g., deploying distributed tracing at an intermediate maturity stage) monotonically improves root-cause analysis (RCA). However, in synchronous microservice call graphs, uncalibrated inclusive span latencies propagate downstream delays to upstream callers---a phenomenon we formalize as \textit{Cascading Span Bias}---which can actively degrade diagnostic accuracy below simple CLI baselines if cross-pillar semantic correlation is absent.
\end{enumerate}

Compounding these domain flaws is a pervasive methodological crisis in software engineering maturity research known as the \textbf{Validation Gap}: in a systematic mapping of 237 maturity model studies, Wendler \cite{wendler_2012} demonstrated that $52\%$ of proposed models lack any validation whatsoever, while $97\%$ of validated models rely exclusively on subjective qualitative interviews or self-reported Likert questionnaires rather than objective, real-system empirical evaluation (Type-3 evaluation in Helgesson et al.'s taxonomy \cite{helgesson_2012}). Similarly, recent academic attempts to define observability maturity (e.g., Patel \cite{patel_2025}) present coarse 3-tier conceptual sketches devoid of mathematical formalization, assessment instruments, or experimental verification.

To bridge this gap, this paper develops, formalizes, and empirically validates the \textbf{Cloud-Native Observability Maturity Model (CN-OMM)}. Strictly adhering to Becker et al.'s eight mandatory DSR requirements (R1--R8) \cite{becker_2009}, De Bruin et al.'s stage-gate design methodology \cite{debruin_2005}, and Lasrado et al.'s configurational theory \cite{lasrado_2015}, we address three research questions:
\begin{itemize}[leftmargin=*, nosep]
    \item \textbf{RQ1 (Structural \& Theoretical Formalization):} \textit{How can cloud-native observability maturity be formally modeled across orthogonal architectural dimensions to account for the dialectical trade-off between diagnostic state reconstructibility and instrumentation probe overhead?}
    \item \textbf{RQ2 (Automated Declarative Assessment):} \textit{How can subjective questionnaire-based maturity assessments be replaced by an objective, non-compensatory mathematical scoring protocol auditable directly against Kubernetes and OpenTelemetry cluster metadata?}
    \item \textbf{RQ3 (Type-3 Empirical Chaos Validation):} \textit{To what extent does progressing across CN-OMM levels yield statistically significant improvements in automated root-cause localization ($MTTDiag, Acc$) under controlled Chaos Engineering fault injection within a bounded CPU overhead envelope?}
\end{itemize}

The primary contributions of this work are fourfold:
\begin{enumerate}[leftmargin=*, nosep]
    \item \textbf{Mathematical \& Configurational Grounding:} We bridge Kalman's control-theoretic state observability \cite{kalman_1960} with distributed telemetry tuples $\mathcal{O}(t) = \langle M, L, T, E, P \rangle$, formalizing maturity progression through Teleological and Dialectical process theories \cite{vandeven_1995} and non-compensatory set-theoretic stage-gates \cite{fiss_2011, lasrado_2015}.
    \item \textbf{The $5 \times 5$ CN-OMM Matrix:} We define 5 evolutionary maturity levels (*Opaque, Reactive, Fragmented, Correlated, Autonomous*) intersected by 5 orthogonal dimensions and 15 sub-components, explicitly incorporating *Resilience Verification via Chaos Engineering* ($D_4$) and *Instrumentation Economics \& Probe Effect Governance* ($D_5$).
    \item \textbf{Automated Maturity Assessment Protocol:} Following Proen\c{c}a's call for automated maturity verification \cite{proenca_2016_rcis}, we specify 15 Quantitative Capability Indicators (QCIs) and an open-source assessment engine that computes effective stage-gate maturity and Euclidean progression indices directly from cluster telemetry.
    \item \textbf{Type-3 Real-System Empirical Validation:} We execute a controlled $2\times2$ factorial Chaos Engineering campaign ($N=30$ paired microservice trials plus $N=60$ chaos injector verification trials) on Kubernetes, empirically proving the collapse of monomodal Level 3 tracing ($0.00\%$ accuracy) and the superiority of Level 4 multimodal correlation ($83.33\%$ accuracy, $17.27\,\text{s}$ median $MTTDiag$, $9.51\%$ CPU overhead).
\end{enumerate}

The remainder of this paper is organized as follows. Section~\ref{sec:background} establishes the control-theoretic and DSR meta-theoretical foundations. Section~\ref{sec:critical_analysis} audits existing academic and industrial observability models (Requirement R1). Section~\ref{sec:cn_omm} presents the mathematical and architectural formalization of the CN-OMM. Section~\ref{sec:assessment} details the automated assessment protocol and scoring equations. Section~\ref{sec:empirical_validation} reports the Type-1 and Type-3 empirical chaos validation. Section~\ref{sec:discussion} discusses architectural implications and threats to validity, and Section~\ref{sec:conclusion} concludes the paper.
